\subsection{氢样离子、约化质量与能级极限}
只有一个电子、核电荷为 $+Ze$ 的氢样离子，可以把库仑常数因子中的 $e^2$ 换成 $Ze^2$。若核质量为 $M$，电子和核均绕质心运动，相对坐标的质量为约化质量
\[
 \mu=\frac{m_eM}{m_e+M}.
\]
相应的轨道半径（两粒子的间距）和能级为
\[
 r_n=\frac{m_e}{\mu}\frac{a_0n^2}{Z},\qquad
 E_n=-\frac{\mu}{m_e}\frac{13.6Z^2}{n^2}\,\mathrm{eV}.
\]
因此氢样氦离子 $\mathrm{He}^+$ 的基态电离能约为 $54.4\,\mathrm{eV}$；中性氦有两个电子，不能直接套这条单电子公式。对同一 $Z$ 的同位素，较重的核使 $\mu$ 更接近 $m_e$，同一跃迁的光子能量略增、波长略短。
这里 $13.6\,\mathrm{eV}$ 已取近似值，精密同位素位移还需一致使用常数精度及其他修正。玻尔近似的速度尺度为 $Z\alpha c/n$；当它不再远小于 $c$ 时，非相对论轨道模型不适用。
